\clearpage
\appendix
\section{Supplementary material}
\subsection{Prompt given to \astra}
System message:
\begin{lstlisting}
You are an engineering drafting and analysis assistant. Use your own reasoning; no external tools are provided. Do not claim you ran a solver.
\end{lstlisting}
User message for a circuit task, verbatim up to the source SVG, which follows in full:
\begin{lstlisting}
Instruction: Set R3 to 150 Ω.

Drawing conventions: straight lines are ideal wires. A wire that ends on another wire (a T) is connected to it; a filled dot also marks a connection; wires that cross without a dot are NOT connected. Rectangles are resistors and circles are ideal DC voltage sources; the '+' inside a source marks its positive terminal. Labels give component names and values.
Criteria (from the drawing notes): every resistor's dissipated power must not exceed the stated rating, and the magnitude of every source's current must not exceed the stated supply fuse rating.

Tasks:
1. Apply the instruction to the drawing, changing nothing else, and return the complete edited SVG.
2. Analyse the EDITED design and report whether it meets every criterion.

Answer with exactly two fenced blocks and nothing else: first ```svg containing the full edited SVG, then ```json containing {"verdict": "pass" or "fail", "violations": [names of every resistor or source that violates a criterion], "source_currents_A": {"V1": current delivered out of its + terminal, ...}, "max_resistor_power_W": {"name": "R?", "value": watts}}

Source SVG:
[complete source SVG]
\end{lstlisting}
For pipe networks the conventions paragraph and the JSON schema are replaced by:
\begin{lstlisting}
Drawing conventions: heavy lines are pipes; filled dots are junctions. A pipe runs between junctions or reservoirs and may turn through 90-degree elbows; each elbow adds the minor-loss coefficient K given in the notes. 'P# L m ØD' gives a pipe's length (m) and internal diameter (mm); 'J# q=.. L/s z=.. m' gives the demand withdrawn at a junction and its elevation; 'R# H=.. m' is a reservoir's fixed total head.
Analysis model: steady incompressible water, density 998 kg/m3, dynamic viscosity 1.002e-3 Pa s, g = 9.80665 m/s2; Darcy-Weisbach head loss with the Churchill (1977) friction factor and the pipe roughness in the notes; tee and junction losses neglected; pressure head at a junction = hydraulic head minus elevation.
Criteria (from the drawing notes): every pipe velocity must not exceed the maximum velocity, and every junction pressure head must be at least the minimum pressure head.

JSON schema: {"verdict": "pass" or "fail", "violations": [names of every pipe or junction that violates a criterion], "pipe_velocities_m_s": {"P1": speed, ...}, "min_pressure_head_m": {"name": "J?", "value": metres}}
\end{lstlisting}

\subsection{Results by edit type}
\Cref{tab:kinds} breaks end-to-end accuracy down by edit type. The one type on which \astra is below 100\,\%
is the pipe-length change that produced its silent wrong verdict (\cref{fig:failure}). Positional addressing
fails mostly on circuit edits, and id addressing is correct on every main-tier type at both model sizes.
\begin{table*}[t]
\centering
\footnotesize
\setlength{\tabcolsep}{5pt}
\begin{tabular}{@{}lcccc@{}}
\toprule
Edit type & \astra & 1.5B pos. & 1.5B id & 7B id \\
\midrule
\kindRows
\bottomrule
\end{tabular}
\caption{End-to-end accuracy (edit, verdict and violations all correct), \% (count). \astra on its completed main-tier tasks; \ours editors (positional or id addressing) on the full 1{,}000-edit main test set.}
\label{tab:kinds}
\end{table*}

\subsection{Training stability}
The first 7B training run, on an H100 NVL, diverged: the loss rose from 0.65 at step 20 to $3.0{\times}10^{4}$
at step 30 with a non-finite gradient norm, and the resulting adapter emits only ``\texttt{!}'' characters. Two
50-step reruns with identical code, settings and data order isolate the cause to PyTorch's cuDNN kernel for
scaled dot-product attention. With it enabled the gradient norm again became non-finite at step 30. With it disabled
all 50 steps were finite and the loss fell from 0.55 to 0.026. The reported 7B editor was retrained with that kernel
disabled. The 1.5B editors, trained on an RTX PRO 6000, never used it. The diverged run is retained,
marked void, and excluded from every table, and the trainer now stops at the first non-finite loss or gradient.
